% !TeX encoding = UTF-8
\documentclass[variant=apuntes,cover=institutional,mode=screen]{ua-engineering-v6-article}
% !TeX encoding = UTF-8
\UASetUniversity{Universidad Austral}
\UASetFaculty{Facultad de Ingeniería}
\UASetDepartment{Departamento de Computación}
\UASetCampus{Pilar}
\UASetCareer{Ingeniería Informática}
\UASetCommission{Comisión A}
\UASetProfessor{Equipo docente}
\UASetSemester{Segundo cuatrimestre 2026}
\UASetCourse{Sistemas Distribuidos}
\UASetDate{2026}


\UASetClassNumber{02}
\UASetClassTitle{RPC, timeouts, retries e idempotencia}
\UASetDocKind{Apuntes}

\title{Clase 02 --- RPC, timeouts, retries e idempotencia}
\author{\UAProfessor}
\date{\UADate}

\begin{document}
\maketitle

\section{Introducción}

La clase 1 mostró que el problema central de los sistemas distribuidos es la \textbf{indistinguibilidad observacional}. La clase 2 estudia cómo esa limitación impacta directamente en el diseño de llamadas remotas.

El objetivo es entender que:

\begin{uakeybox}
Una llamada remota no es una función distribuida: es una interacción bajo incertidumbre.
\end{uakeybox}

\section{RPC: abstracción y límites}

RPC (Remote Procedure Call) busca abstraer una invocación remota como si fuera una llamada local. Sin embargo, esta abstracción es incompleta.

\subsection{Diferencias fundamentales}

\begin{itemize}
\item una llamada local no falla parcialmente;
\item una llamada remota sí;
\item una llamada local tiene semántica clara;
\item una llamada remota depende de la red.
\end{itemize}

\section{Modelo de ejecución}

Sea una operación remota $op$.

Eventos posibles:

\begin{itemize}
\item envío del request
\item recepción en servidor
\item ejecución
\item envío de respuesta
\item recepción en cliente
\end{itemize}

Cada etapa puede fallar independientemente.

\section{No-inmediatez y capas de comunicación}

Una comunicación distribuida no debe entenderse como un salto instantáneo entre cliente y servidor. Entre la intención de invocar una operación y la observación de una respuesta hay una trayectoria compuesta por capas: aplicación, librería RPC o cliente HTTP, serialización, transporte, red, balanceadores, servidor, almacenamiento y camino de regreso.

Cada capa puede reaccionar de manera distinta ante una falla:
\begin{itemize}
\item la aplicación puede abandonar la espera y mostrar error;
\item el cliente RPC puede reintentar automáticamente;
\item TCP puede retransmitir segmentos sin que la aplicación lo observe directamente;
\item un balanceador puede redirigir a otra instancia;
\item el servidor puede seguir procesando aunque el cliente ya haya vencido su timeout;
\item la base de datos puede confirmar o abortar una operación independientemente de que la respuesta vuelva.
\end{itemize}

Esto introduce una forma práctica de concurrencia: mientras una capa decide que ``ya esperó suficiente'', otra capa puede seguir intentando, entregando, ejecutando o confirmando efectos. Para diagnosticar un problema no alcanza con preguntar ``falló o no falló'': hay que reconstruir qué observó cada capa, qué acción tomó y qué efecto pudo producir.

\begin{uakeybox}
La no-inmediatez de la comunicación obliga a razonar con líneas de tiempo: lo que el cliente observa, lo que la red transporta, lo que el servidor ejecuta y lo que el almacenamiento confirma pueden no coincidir en el mismo instante.
\end{uakeybox}

\section{Semánticas}

\subsection{At-least-once}

El sistema reintenta hasta obtener respuesta.

Ventaja:
\begin{itemize}
\item mejora liveness
\end{itemize}

Problema:
\begin{itemize}
\item múltiples ejecuciones
\end{itemize}

\subsection{At-most-once}

Evita duplicación.

Problema:
\begin{itemize}
\item puede perder ejecución
\end{itemize}

\subsection{Exactly-once}

\begin{uakeybox}
En sistemas distribuidos reales, exactly-once requiere coordinación fuerte o se implementa como combinación de at-least-once + idempotencia.
\end{uakeybox}

\section{Timeouts}

Un timeout define cuándo el cliente deja de esperar.

\subsection{Propiedades}

\begin{itemize}
\item decisión local
\item basada en tiempo
\item introduce incertidumbre
\end{itemize}

\section{Retries}

Retry = reintento tras timeout.

\subsection{Problema estructural}

Retry introduce duplicación potencial porque:

\begin{itemize}
\item no se sabe si la ejecución ocurrió
\item se reintenta igual
\end{itemize}

\section{Idempotencia}

Una operación es idempotente si múltiples ejecuciones producen el mismo estado final.

\subsection{Formalización}

\[
op(op(s)) = op(s)
\]

\subsection{Implementación}

\begin{itemize}
\item request ID
\item log de operaciones
\item cache de resultados
\end{itemize}

\section{Trade-offs}

\begin{itemize}
\item idempotencia → más estado
\item retries → más robustez
\item timeouts → más incertidumbre
\end{itemize}

\section{Conclusión}

\begin{uakeybox}
El diseño de RPC consiste en elegir explícitamente:
\begin{itemize}
\item semántica de ejecución
\item política de timeout
\item política de retry
\item mecanismo de idempotencia
\end{itemize}
\end{uakeybox}

\end{document}
